\title{Rebuttal Y/2}

\maketitle

\section{Comments}

\begin{itemize}
    \item Keep the bugs internalized in their respective code files and keep a general summary in the report
    \item Keep the report small (3 pages at most)
    \item Something very wrong with the partial view point-cloud. Try: first sample points by, given the camera pose, projecting the pixels sampled on the object silhouette (or something along the lines of that) and get the 3D ray-intersection position from the intersecting geometry to actually get a proper partial view point cloud? I think the last iteration in the `/OLD` directory has this code properly implemented somewhere.
    \item The frustum looks atrocious and wrong. The fustum might be too big in the visualization, try smallar frustum size and make sure it is correct.
    \item Would be great if the visualization is in GIF to spatially show the visualization correctly without needing to open up an expensive local 3D viewer to self-verify (you and me.)
    \item Add re-decoded GT mesh rendering (GIF?) with the output latents explicitly to self-verify we are correctly using TRELLIS2's VAE.
\end{itemize}

\section{Verification of the Training Pipeline}

Once the aforementioned lists are done, train an end-to-end pipeline on ONE OBJECT with MANY PARTIAL VIEWS. To check the training is working correctly, we use the trained weights to run inference on the \textit{same} object used to train but with \textit{different} camera pose (i.e., view.). Explicitly reference the weight checkpoints and make sure wandb logging is done. I want result, you want result, let's get it verified.

